\documentclass[twoside,10pt]{article}

\usepackage[T1]{fontenc}
\usepackage{lmodern}
\usepackage{amsmath,amsthm,amssymb,amsfonts}
\usepackage{mathtools}
\usepackage[
  letterpaper,
  total={5.5in,8in},
  vmarginratio=1:1,
  hmarginratio=1:1,
  headsep=21pt
]{geometry}
\usepackage{xcolor}
\usepackage[
  breaklinks=false,
  linktocpage=true
]{hyperref}
\hypersetup{
  colorlinks=true,
  linkcolor={purple},
  urlcolor={purple},
  citecolor={blue!80!black},
  pdfauthor={},
  pdftitle={A Single Kinetic Defect in a Run-and-Tumble Ring:
    Exact Spectrum, Exceptional Points, and Relaxation Crossover},
  pdfsubject={Exact spectral analysis of a run-and-tumble ring with one tumbling-rate defect},
  pdfkeywords={run-and-tumble particle, kinetic defect, exceptional point,
    localized mode, spectral gap, telegraph equation}
}

\medmuskip=4mu
\thickmuskip=8mu
\setlength{\emergencystretch}{2em}

\numberwithin{equation}{section}

\newtheoremstyle{slbody}
  {8pt}
  {8pt}
  {\normalfont\slshape}
  {}
  {\bfseries}
  {.}
  {0.5em}
  {\thmname{#1}\thmnumber{ #2}%
   \thmnote{ {\normalfont\itshape(\ignorespaces#3\unskip)}}}
\theoremstyle{slbody}
\newtheorem{theorem}{Theorem}[section]
\newtheorem{proposition}[theorem]{Proposition}
\newtheorem{lemma}[theorem]{Lemma}
\newtheorem{corollary}[theorem]{Corollary}
\theoremstyle{definition}
\newtheorem{definition}[theorem]{Definition}
\theoremstyle{remark}
\newtheorem{remark}[theorem]{Remark}

\renewcommand{\qedsymbol}{$\blacksquare$}

\newcommand{\Z}{\mathbb Z}
\newcommand{\C}{\mathbb C}
\newcommand{\one}{\mathbf 1}
\newcommand{\ii}{\mathrm i}
\newcommand{\spec}{\operatorname{spec}}
\newcommand{\ord}{\operatorname{ord}}
\newcommand{\Res}{\operatorname*{Res}}

\title{A single kinetic defect in a run-and-tumble ring:\\
exact spectrum, exceptional points, and relaxation crossover}
\author{}
\date{}

\makeatletter
\renewcommand{\section}{%
  \@startsection{section}{1}{\z@}%
    {-3.5ex \@plus -1ex \@minus -.2ex}%
    {2.3ex \@plus .2ex}%
    {\normalfont\large\bfseries}}
\renewcommand{\@seccntformat}[1]{%
  \csname the#1\endcsname.\enspace}

\def\tagform@#1{%
  \maketag@@@{(\ignorespaces{\oldstylenums{#1}}\unskip\@@italiccorr)}}
\renewcommand{\eqref}[1]{%
  \textup{{\normalfont(\oldstylenums{\ref{#1}}}\normalfont)}}

\newcommand{\shorttitle}{A SINGLE KINETIC DEFECT IN A RUN-AND-TUMBLE RING}
\newcommand{\shortauthors}{}
\newcommand{\ps@bookheader}{%
  \renewcommand\@oddfoot{\hfil}%
  \renewcommand\@evenfoot{\hfil}%
  \renewcommand\@oddhead{%
    \ifnum\value{page}>1
      {\small\hfil\shortauthors}\hfil\thepage
    \else
      \hfil
    \fi}%
  \renewcommand\@evenhead{%
    \thepage\hfil\small\shorttitle\hfil}}
\pagestyle{bookheader}

\renewcommand{\maketitle}{%
  \begin{center}
    {\large\bfseries\@title}\\
    \vskip36pt
    \ifx\@author\@empty
      \vspace{55pt}\par
    \else
      {\scshape\@author}\par
      \medskip
      {\small\@date}\par
    \fi
    \vskip36pt
  \end{center}}
\makeatother

\renewenvironment{abstract}
  {\quotation\noindent\small{\bfseries\abstractname.\enspace}}
  {\endquotation}

\begin{document}
\maketitle

\begin{abstract}
We solve the spectral problem for one continuous-time run-and-tumble
particle on an $n$-site ring when the tumbling rate at one site is
$\omega _0$ and the bulk rate is $\omega$.  The defect changes the
$2n\times2n$ generator by rank one.  The characteristic polynomial is
the homogeneous polynomial plus $(\omega _0-\omega)/n$ times its
$\omega$-derivative.  This identity gives Chebyshev and transfer-matrix
quantization formulas, a parity factorization, and finite polynomial
tests for every Jordan block.  We separately resolve the singular
flat-band point $\omega=1$, where the eigenvalue $-2$ can carry Jordan
blocks of sizes two or three.  The infinite-line Green function gives
an exact criterion for exponentially localized defect modes.  For fixed
positive bulk rate, reflection protects one homogeneous long-wave mode;
consequently a slow defect, including $\omega _0\asymp n^{-1}$, does not
produce a new spectral-gap scaling.  The nontrivial $n^{-1}$ crossover
occurs when the bulk rate itself is of order $n^{-1}$.  Finally, the
standard telegraph equation misses the lattice Poisson-diffusion term at
fixed $\omega$, while a diffusion-corrected telegraph equation agrees
through order $n^{-2}$.
\end{abstract}

\vskip5pt
\noindent\textbf{Keywords.}\enspace Run-and-tumble particle; kinetic
defect; non-Hermitian spectrum; exceptional point; localized mode;
spectral gap; telegraph equation.

\vskip50pt
\baselineskip=13.5pt

\section{Introduction}

Persistent random walks provide a microscopic description of motion with
a finite directional memory.  Continuum formulations with
space-dependent turning statistics arise, for example, in the theory of
chemotactic random walks \cite{Schnitzer1993}, while lattice
run-and-tumble models retain the individual hopping and reversal events
and connect directly to nonequilibrium particle systems
\cite{ThompsonTailleurCatesBlythe2011}.  For a homogeneous ring, the
one-particle spectrum, its exceptional points, and the associated
relaxation transition are known explicitly
\cite{MallminBlytheEvans2019}.
Spatial inhomogeneity has been studied through position-dependent speed
\cite{AngelaniGarra2019}, quenched disorder in speed and tumbling
\cite[Secs.~2.1.2--2.1.3 and 2.2.2--2.2.3]
{BenDorWoillezKafriKardarSolon2019}, and saturating space-dependent rates
on the line \cite[Sec.~II, Eqs.~(3)--(4)]{JainChatterjee2024}.  Those
works emphasize probability laws, steady states, or relaxation in
extended media.  The present problem instead asks for the full
finite-ring characteristic polynomial and Jordan structure of a single
lattice tumbling defect, together with its large-ring spectral ordering.

This paper asks how that spectral picture changes when the tumbling rate
is altered at one site.  The inhomogeneity is local but the generator is
non-normal, so eigenvalue crossings, Jordan structure, and long-wave
relaxation must be tracked together.  The central observation is that
the defect is rank one in the polarization channel.  It leads to a
finite characteristic identity, a parity decomposition, an
infinite-line Green-function criterion including a compact defect mode,
and a direct comparison between lattice and telegraph relaxation.  The
fixed-rate and bulk-critical limits separate the effects of a local
defect from those of a vanishing bulk tumbling rate.

\section{Model and notation}

\begin{definition}[Defective run-and-tumble generator]\label{def:generator}
Let $n\geq2$, let $\omega>0$, and let $\omega _0\geq0$.  Sites are
$j\in\Z/n\Z$, orientations are $\sigma\in\{+,-\}$, and
\[
  \omega_j=\begin{cases}\omega_0,&j=0,\\ \omega,&j\ne0.\end{cases}
\]
The forward generator $L=L_n(\omega,\omega_0)$ is
\begin{align}
 (Lu)^+_j&=u^+_{j-1}-(1+\omega_j)u^+_j+\omega_j u^-_j,\label{eq:generator-plus}\\
 (Lu)^-_j&=u^-_{j+1}-(1+\omega_j)u^-_j+\omega_j u^+_j.\label{eq:generator-minus}
\end{align}
Thus hopping in the current orientation has rate one and tumbling has
rate $\omega_j$.  Put
\[
 \delta=\omega_0-\omega,\qquad
 q_k=\frac{2\pi k}{n},\qquad c_k=\cos q_k,\qquad s_k=\sin q_k.
\]
\end{definition}

Equations \eqref{eq:generator-plus}--\eqref{eq:generator-minus} use the
same one-particle convention as equations (6)--(10) of
\cite{MallminBlytheEvans2019}.  That paper records the homogeneous
eigenvalues in its equation (10), the homogeneous exceptional points in
equation (13), and the first-mode relaxation rate in equations
(20)--(21).

\begin{lemma}[Rank-one defect]\label{lem:rank-one}
Let $L^{(0)}=L_n(\omega,\omega)$.  If
$d=e_{0,+}-e_{0,-}$, then
\[
 L=L^{(0)}-\delta\,dd^{\mathsf T}.
\]
In particular, the perturbation has rank one when $\delta\ne0$.
\end{lemma}

\begin{proof}
At site zero, changing $\omega$ to $\omega_0$ adds
\[
 \delta\begin{pmatrix}-1&1\\1&-1\end{pmatrix}
 =-\delta\begin{pmatrix}1\\-1\end{pmatrix}
             \begin{pmatrix}1&-1\end{pmatrix}.
\]
Every other matrix entry is unchanged.
\end{proof}

\section{Three exact characteristic equations}

\begin{definition}[Homogeneous factors]\label{def:homogeneous-factors}
For $\lambda\in\C$, set
\begin{align}
 a&=\lambda+1+\omega,\label{eq:a-def}\\
 x&=\frac{a^2+1-\omega^2}{2a},\label{eq:x-def}\\
 D_k(\lambda)&=(\lambda+1-c_k)(\lambda+1+2\omega-c_k)+s_k^2
             =a^2+1-\omega^2-2ac_k.\label{eq:Dk}
\end{align}
Although $x$ is undefined at $a=0$, each expression below containing
$a^nT_n(x)$ is understood by its unique polynomial continuation.
\end{definition}

\begin{lemma}[Homogeneous polynomial]\label{lem:homogeneous-polynomial}
The homogeneous characteristic polynomial is
\begin{align}
 P_n^{(0)}(\lambda;\omega)
  &=\det(\lambda I-L^{(0)})\notag\\
  &=\prod_{k=0}^{n-1}D_k(\lambda)\label{eq:P0-product}\\
  &=2a^n\bigl(T_n(x)-1\bigr),\label{eq:P0-cheb}
\end{align}
where $T_n$ is the Chebyshev polynomial of the first kind.
The two roots of $D_k$ are
\[
 \lambda_k^\pm=c_k-1-\omega
       \pm\sqrt{\omega^2-s_k^2}.
\]
\end{lemma}

\begin{proof}
Fourier transformation gives the block
\[
 W_k=\begin{pmatrix}e^{\ii q_k}-1-\omega&\omega\\
                     \omega&e^{-\ii q_k}-1-\omega\end{pmatrix}.
\]
Its characteristic polynomial is $D_k$, which proves
\eqref{eq:P0-product} and the displayed roots.  Since
$D_k=2a(x-c_k)$ and
\[
 \prod_{k=0}^{n-1}(x-c_k)=2^{1-n}\bigl(T_n(x)-1\bigr),
\]
equation \eqref{eq:P0-cheb} follows.
\end{proof}

\begin{theorem}[Exact finite characteristic polynomial]\label{thm:characteristic}
For every $n\geq2$, $\omega>0$, and $\omega_0\geq0$,
\begin{align}
 P_n(\lambda;\omega,\omega_0)
  &:=\det(\lambda I-L_n(\omega,\omega_0))\label{eq:P-def}\\
  &=P_n^{(0)}(\lambda;\omega)
    +\frac{\delta}{n}\,\partial_\omega
       P_n^{(0)}(\lambda;\omega)\label{eq:derivative-form}\\
  &=P_n^{(0)}(\lambda;\omega)
    \left[1+\frac{2\delta}{n}
      \sum_{k=0}^{n-1}\frac{\lambda+1-c_k}{D_k(\lambda)}\right].
      \label{eq:resolvent-form}
\end{align}
Equation \eqref{eq:resolvent-form} is interpreted by polynomial
continuation at a homogeneous eigenvalue.  Equivalently,
\begin{align}
 P_n={}&2a^n(T_n(x)-1)
      +2\delta a^{n-1}(T_n(x)-1)
      +\delta a^{n-2}\lambda(\lambda+2)U_{n-1}(x),
      \label{eq:cheb-defect}
\end{align}
where $U_{n-1}$ is the Chebyshev polynomial of the second kind and the
right-hand side again denotes its polynomial continuation.
\end{theorem}

\begin{proof}
Let $B=\lambda I-L^{(0)}$.  Lemma \ref{lem:rank-one} and the matrix
determinant lemma give
\[
 \det(\lambda I-L)=\det B\bigl(1+\delta d^{\mathsf T}B^{-1}d\bigr).
\]
In Fourier space,
\[
 d^{\mathsf T}B^{-1}d
   =\frac1n\sum_{k=0}^{n-1}
      \frac{2(\lambda+1-c_k)}{D_k(\lambda)},
\]
which proves \eqref{eq:resolvent-form}.  Direct differentiation of
\eqref{eq:P0-product} gives
\[
 \partial_\omega P_n^{(0)}
 =P_n^{(0)}\sum_k\frac{2(\lambda+1-c_k)}{D_k},
\]
and hence \eqref{eq:derivative-form}.  Finally,
\[
 \partial_\omega x=\frac{\lambda(\lambda+2)}{2a^2},
 \qquad T_n'(x)=nU_{n-1}(x),
\]
and differentiating \eqref{eq:P0-cheb} proves
\eqref{eq:cheb-defect}.
\end{proof}

\begin{theorem}[Transfer-matrix quantization]\label{thm:transfer}
For a local rate $r\geq0$, define
\[
 T_r(\lambda)=\frac1{\lambda+1+r}
 \begin{pmatrix}
  1&r\\
  -r&(\lambda+1+r)^2-r^2
 \end{pmatrix}.
 \label{eq:transfer}
\]
Then $\det T_r=1$, and, away from the displayed denominators,
\[
 P_n(\lambda;\omega,\omega_0)
 =a^{n-1}(\lambda+1+\omega_0)
   \left\{\operatorname{tr}
     \bigl[T_\omega(\lambda)^{n-1}T_{\omega_0}(\lambda)\bigr]-2\right\}.
 \label{eq:transfer-char}
\]
After multiplication by the denominators, this identity holds for every
$\lambda$ by polynomial continuation.  When $a\ne0$ and
$\lambda+1+\omega_0\ne0$, an eigenvalue is equivalently characterized by
\[
 \operatorname{tr}
     \bigl[T_\omega(\lambda)^{n-1}T_{\omega_0}(\lambda)\bigr]=2.
 \label{eq:trace-two}
\]
At a zero of either denominator, the cleared polynomial identity
\eqref{eq:transfer-char}, rather than the undefined trace in
\eqref{eq:trace-two}, is the eigenvalue criterion.
Moreover, if $t=\operatorname{tr}T_\omega=2x$, then
\[
 T_\omega^{n-1}=U_{n-2}(x)T_\omega-U_{n-3}(x)I,
 \label{eq:transfer-cheb}
\]
with $U_{-1}=0$.
\end{theorem}

\begin{proof}
The eigenvalue equations at site $j$ are
\begin{align*}
 (\lambda+1+\omega_j)u_j^+-\omega_j u_j^-&=u_{j-1}^+,\\
 (\lambda+1+\omega_j)u_j^--\omega_j u_j^+&=u_{j+1}^-.
\end{align*}
Consequently
\[
 \binom{u_j^+}{u_{j+1}^-}
 =T_{\omega_j}(\lambda)\binom{u_{j-1}^+}{u_j^-}.
\]
Periodic boundary conditions require the product monodromy to have
eigenvalue one.  Its determinant is one, so this is equivalent to trace
two.  To identify the polynomial, retain temporarily arbitrary rates
$r_0,\ldots,r_{n-1}$, denote the resulting generator by
$L_n(r_0,\ldots,r_{n-1})$, and put $b_j=\lambda+1+r_j$.  Gaussian elimination
of the $2n$ eigenvalue equations, in the order
$u_0^+,u_1^+,\ldots,u_{n-1}^+$, has pivots $b_0,\ldots,b_{n-1}$.
If $M=T_{r_{n-1}}\cdots T_{r_0}$, the two remaining cyclic boundary
equations have coefficient matrix $M-I$.  The cyclic row ordering
contributes one minus sign, and
$-\det(M-I)=\operatorname{tr}M-2$ because $\det M=1$.  Consequently
\begin{equation}\label{eq:general-transfer-determinant}
 \det(\lambda I-L_n(r_0,\ldots,r_{n-1}))
 =\left(\prod_{j=0}^{n-1}b_j\right)
  \left[\operatorname{tr}(T_{r_{n-1}}\cdots T_{r_0})-2\right].
\end{equation}
This is an identity of rational functions where all $b_j\ne0$;
multiplication by the displayed pivots makes both sides polynomials, so
the identity extends to every $\lambda$.  Taking $r_0=\omega_0$ and all
other rates equal to $\omega$, and using cyclicity of the trace, gives
\eqref{eq:transfer-char}, including algebraic multiplicities.  Equation
\eqref{eq:transfer-cheb} is the Cayley--Hamilton identity for a
unimodular $2\times2$ matrix.
\end{proof}

\section{Parity factorization and every finite exceptional point}

\begin{definition}[Kinetic reflection]\label{def:reflection}
Define the involution
\[
 (Ru)^+_j=u^-_{-j},\qquad (Ru)^-_j=u^+_{-j}.
\]
Write $\mathcal H_\pm=\ker(R\mp I)$.
\end{definition}

\begin{lemma}[Protected sector]\label{lem:protected-sector}
The generator commutes with $R$, the defect vector satisfies $Rd=-d$,
and
\[
 L|_{\mathcal H_+}=L^{(0)}|_{\mathcal H_+}.
\]
Thus the defect changes only the reflection-odd sector.
\end{lemma}

\begin{proof}
Reflection reverses position and orientation simultaneously, leaving the
directed hopping rules and the rate field $\omega_j$ invariant.  At the
fixed site zero it exchanges the two entries of $d$, hence $Rd=-d$.
Lemma \ref{lem:rank-one} then vanishes on $\mathcal H_+$.
\end{proof}

\begin{definition}[Parity polynomials]\label{def:parity-polynomials}
Put $m=\lfloor(n-1)/2\rfloor$ and let $\varepsilon_n$ be one for even
$n$ and zero for odd $n$.  Define
\begin{align}
 H_{n,+}(\lambda)
   &=\lambda(\lambda+2)^{\varepsilon_n}
     \prod_{k=1}^{m}D_k(\lambda),\label{eq:Hplus}\\
 H_{n,-}^{(0)}(\lambda)
   &=(\lambda+2\omega)
     (\lambda+2+2\omega)^{\varepsilon_n}
     \prod_{k=1}^{m}D_k(\lambda),\label{eq:Hminus0}\\
 G_n(\lambda)
   &=\frac1n\left[
       \frac1{\lambda+2\omega}
       +\frac{\varepsilon_n}{\lambda+2+2\omega}
       +2\sum_{k=1}^{m}\frac{\lambda+1-c_k}{D_k(\lambda)}
     \right],\label{eq:Gfinite}\\
 Q_{n,-}(\lambda)
   &=H_{n,-}^{(0)}(\lambda)\bigl(1+2\delta G_n(\lambda)\bigr).
   \label{eq:Qminus}
\end{align}
Every quotient in \eqref{eq:Qminus} cancels, so $Q_{n,-}$ is a monic
polynomial of degree $n$.  An entirely polynomial version is
\begin{align}
Q_{n,-}={}&H_{n,-}^{(0)}+\frac{2\delta}{n}\left[
 \frac{H_{n,-}^{(0)}}{\lambda+2\omega}
 +\varepsilon_n\frac{H_{n,-}^{(0)}}{\lambda+2+2\omega}\right.\notag\\
&\left.\hspace{38mm}
 +2\sum_{k=1}^{m}(\lambda+1-c_k)
       \frac{H_{n,-}^{(0)}}{D_k}\right].\label{eq:Qminus-poly}
\end{align}
\end{definition}

\begin{theorem}[Exact parity factorization]\label{thm:parity-factor}
The characteristic polynomial factors as
\[
 P_n(\lambda;\omega,\omega_0)
   =H_{n,+}(\lambda)Q_{n,-}(\lambda).
 \label{eq:parity-factor}
\]
The first factor is the characteristic polynomial on $\mathcal H_+$;
the second is the characteristic polynomial on $\mathcal H_-$.
\end{theorem}

\begin{proof}
The modes $q_k$ and $-q_k$ form a four-dimensional Fourier subspace,
which splits into two isomorphic two-dimensional $R$-eigenspaces.  Each
has characteristic polynomial $D_k$.  At $q=0$, the vectors $(1,1)$ and
$(1,-1)$ have parities $+$ and $-$ and eigenvalues $0$ and $-2\omega$.
When $n$ is even, the $q=\pi$ density and polarization vectors have
parities $+$ and $-$ and eigenvalues $-2$ and $-2-2\omega$.  This gives
\eqref{eq:Hplus}--\eqref{eq:Hminus0}.  Lemma
\ref{lem:protected-sector} and the rank-one determinant lemma restricted
to $\mathcal H_-$ give \eqref{eq:Qminus}.  Multiplication proves
\eqref{eq:parity-factor}.
\end{proof}

\begin{definition}[Exceptional point]\label{def:exceptional}
For an eigenvalue $\mu$, let $m_{\rm alg}(\mu)$ be its algebraic
multiplicity and
$m_{\rm geo}(\mu)=\dim\ker(\mu I-L)$ its geometric multiplicity.  We call
$(\omega,\omega_0,\mu)$ an exceptional point when
$m_{\rm alg}(\mu)>m_{\rm geo}(\mu)$.
Write $J_r(\mu)$ for a Jordan block of size $r$ at $\mu$.
\end{definition}

\begin{lemma}[No inter-wave coincidences away from the flat band]
\label{lem:no-coincidences}
If $c_k\ne c_\ell$ and $D_k(\mu)=D_\ell(\mu)=0$, then $\omega=1$ and
$\mu=-2$.
\end{lemma}

\begin{proof}
Subtracting the two equations gives
$2a(c_k-c_\ell)=0$, hence $a=0$.  Substitution into
\eqref{eq:Dk} gives $1-\omega^2=0$.  Since $\omega>0$, this means
$\omega=1$ and then $\mu=-1-\omega=-2$.
\end{proof}

\begin{theorem}[Jordan classification for $\omega\ne1$]
\label{thm:jordan-regular}
Assume $\omega\ne1$.
\begin{enumerate}
\item In the protected sector, the endpoint factors in
  \eqref{eq:Hplus} give one-dimensional blocks.  For each
  $1\leq k\leq m$, a root of $D_k$ gives one $J_1$, except that
  \begin{equation}\label{eq:protected-ep}
    \omega=|s_k|,\qquad
    \mu=c_k-1-\omega.
  \end{equation}
  gives one $J_2(\mu)$.
\item In the affected sector, if
  \begin{equation}\label{eq:affected-ep}
    Q_{n,-}(\mu)=Q_{n,-}'(\mu)=\cdots
      =Q_{n,-}^{(r-1)}(\mu)=0,\qquad
    Q_{n,-}^{(r)}(\mu)\ne0,
  \end{equation}
  then the complete affected-sector contribution at $\mu$ is one
  block $J_r(\mu)$.
\item The full Jordan form at $\mu$ is the direct sum of the blocks in
  items 1 and 2.  In particular, a crossing between the two parity
  sectors is semisimple unless one of its sector blocks has size greater
  than one.
\end{enumerate}
Consequently, \eqref{eq:protected-ep} together with
\eqref{eq:affected-ep} is an exact finite criterion for every exceptional
point when $\omega\ne1$.  The order $r$ in \eqref{eq:affected-ep}
classifies exceptional points of arbitrary order.
\end{theorem}

\begin{proof}
The discriminant of $D_k$ is $4(\omega^2-s_k^2)$.  At a zero
discriminant its Fourier block is not scalar and hence is a single
$J_2$; otherwise each root is simple.  Lemma
\ref{lem:no-coincidences} excludes coincidences between distinct
$c_k$ values.

It remains to show that the affected restriction is cyclic.  In each
two-dimensional homogeneous parity block, the spectral weight of the
defect vector has numerator $2(\lambda+1-c_k)$.  This numerator and
$D_k$ have no common root for an interior wave number: a common root
would give $\lambda+1=c_k$, after which $D_k=s_k^2\ne0$.  The endpoint
weights in \eqref{eq:Gfinite} are also nonzero.  Lemma
\ref{lem:no-coincidences} makes the block minimal polynomials coprime.
The defect vector is therefore cyclic for the homogeneous odd
restriction.  If $A$ has cyclic vector $d$, then $d$ is also cyclic for
$A-\delta dd^{\mathsf T}$, because each
$(A-\delta dd^{\mathsf T})^jd$ equals $A^jd$ plus a linear combination
of lower powers.  Hence the minimal and characteristic polynomials of
the affected restriction both equal $Q_{n,-}$.  A cyclic matrix has one
Jordan block for each eigenvalue, of size equal to the root order.
The direct-sum assertion follows from the invariant parity splitting.
\end{proof}

\begin{theorem}[The complete flat-band classification at $\omega=1$]
\label{thm:flat-band}
Let $\omega=1$ and put $\mu=-2$.
\begin{enumerate}
\item If $\omega_0=1$ and $4\nmid n$, the contribution at $\mu$ is
  $nJ_1(\mu)$.
\item If $\omega_0=1$ and $4\mid n$, the contribution is
  \[
    2J_2(\mu)\oplus(n-2)J_1(\mu).
  \]
\item If $\omega_0\ne1$ and $n\equiv1$ or $2\pmod4$, the contribution
  is $(n-1)J_1(\mu)$.
\item If $\omega_0\ne1$ and $n\equiv0\pmod4$, the contribution is
  \[
    J_2(\mu)\oplus(n-2)J_1(\mu).
  \]
\item If $n\equiv3\pmod4$, $\omega_0\ne1$, and
  $\omega_0\ne(n-1)/(n+1)$, the contribution is again
  $J_2(\mu)\oplus(n-2)J_1(\mu)$.
\item If $n\equiv3\pmod4$ and
  $\omega_0=(n-1)/(n+1)$, the contribution is
  \[
    J_3(\mu)\oplus(n-2)J_1(\mu).
  \]
\end{enumerate}
For eigenvalues $\lambda\ne-2$, every protected-sector root is a
one-dimensional block.  Remove the exact power of $(\lambda+2)$
specified above from $Q_{n,-}$.  Every root of the remaining affected
polynomial contributes one Jordan block, so the derivative test
\eqref{eq:affected-ep} applied to that reduced polynomial classifies
every remaining affected block at $\omega=1$.  The full
Jordan form is the direct sum of the protected and affected blocks.
\end{theorem}

\begin{proof}
Write $t=\lambda+2$.  At $\omega=1$,
\[
 D_k=t(t-2c_k),\qquad
 P_n=t^{n-1}F_n(t),
\]
where, for $\delta=\omega_0-1$,
\begin{equation}\label{eq:flat-F}
 F_n(t)=2t\bigl(T_n(t/2)-1\bigr)
 +\delta\left[2\bigl(T_n(t/2)-1\bigr)
 +(t-2)U_{n-1}(t/2)\right].
\end{equation}
The values
\[
 T_n(0)=\cos(n\pi/2),\qquad U_{n-1}(0)=\sin(n\pi/2)
\]
give the following algebraic multiplicities.  For $\delta=0$ they are
$n$, except that $4\mid n$ gives $n+2$.  For $\delta\ne0$ they are
$n-1$ when $n\equiv1,2\pmod4$, and $n$ when
$n\equiv0,3\pmod4$.  In the last congruence class,
$F_n'(0)=-2-\delta(n+1)$; it vanishes exactly at
$\delta=-2/(n+1)$.  For $n\equiv3\pmod4$ one has
$U_{n-1}''(0)=n^2-1$, and direct differentiation of
\eqref{eq:flat-F} at the tuned value gives
\begin{equation}\label{eq:flat-second-derivative}
 F_n''(0)=-2n+\frac{U_{n-1}''(0)}{n+1}=-(n+1)\ne0.
\end{equation}
Thus the tuned multiplicity is $n+1$.

The bulk equations at $t=0$ read
\[
 u_j^-=-u_{j-1}^+\qquad(j\ne0).
\]
Together they impose this relation for every $j$.  The defect equation
adds no constraint when $\omega_0=1$ and otherwise adds the single
constraint $u_0^++u_{n-1}^+=0$.  Thus the geometric multiplicity is $n$
in the homogeneous case and $n-1$ otherwise.  Fourier decomposition
shows that the two excess algebraic dimensions when $4\mid n$ and
$\delta=0$ are the two $q=\pm\pi/2$ blocks, each a $J_2$.  For
$\delta\ne0$, the parity factorization places the $4\mid n$ excess in
one protected $J_2$.

It remains to distinguish one $J_3$ from two $J_2$ blocks in the tuned
$n\equiv3\pmod4$ case.  Put $N_0=L_n(1,1)+2I$ and
$N=L_n(1,\omega_0)+2I=N_0-\delta dd^{\mathsf T}$.  Because
$4\nmid n$, zero is a semisimple eigenvalue of $N_0$ with eigenspace
$E=\ker N_0$ of dimension $n$.  Let $\Pi$ be its spectral projection.
Write
\[
 \det(tI-N_0)=t^nh(t),\qquad h(0)\ne0.
\]
Semisimplicity gives
$(tI-N_0)^{-1}=t^{-1}\Pi+O(1)$.  The rank-one determinant formula then
shows that the coefficient of $t^{n-1}$ in $\det(tI-N)$ is
$\delta h(0)d^{\mathsf T}\Pi d$.
The multiplicity computation above shows that this coefficient vanishes
for $n\equiv3\pmod4$, and therefore
\begin{equation}\label{eq:flat-isotropic-projection}
 d^{\mathsf T}\Pi d=0.
\end{equation}
The eigenspace of $N$ is
$K=E\cap\ker d^{\mathsf T}$, which has dimension $n-1$.  If
$v\in K\cap\operatorname{ran}N$, write $v=Nu$ and apply $\Pi$:
\begin{equation}\label{eq:flat-range-line}
 v=\Pi Nu=-\delta\,\Pi d\,(d^{\mathsf T}u).
\end{equation}
Thus $K\cap\operatorname{ran}N$ has dimension at most one.  It has
dimension one because the algebraic multiplicity exceeds the geometric
multiplicity.  Hence exactly one Jordan block at zero for $N$ is
nontrivial.  It is a $J_2$ when the algebraic multiplicity is $n$, and
the additional algebraic order at $\delta=-2/(n+1)$ makes that same
block a $J_3$.  This proves all six cases and excludes the alternative
of two $J_2$ blocks in the tuned case.

For completeness, put
$A_-^{(0)}=L^{(0)}|_{\mathcal H_-}$ and consider an affected eigenvalue
$\lambda\ne-2$.
If $\delta\ne0$, it cannot equal a nonflat homogeneous root: the
corresponding residue in \eqref{eq:Gfinite} is nonzero, so
$Q_{n,-}$ does not vanish at that pole.  Hence
$\lambda I-A_-^{(0)}$ is invertible.  The rank-one eigenvalue equation
makes its eigenspace the one-dimensional span of
$(\lambda I-A_-^{(0)})^{-1}d$.  Its root order therefore gives the size
of its unique Jordan block.  If $\delta=0$, all nonflat roots are simple
by Lemma \ref{lem:no-coincidences}.  This proves the final classification
without asserting cyclicity at the removed flat primary component.
\end{proof}

\begin{corollary}[Finite algebraic exceptional set]\label{cor:resultant}
For fixed $n$, all non-flat affected exceptional parameters lie on the
explicit algebraic set
\[
 \Res_\lambda\bigl(Q_{n,-},\partial_\lambda Q_{n,-}\bigr)=0,
\]
with $Q_{n,-}$ given by \eqref{eq:Qminus-poly}.  At $\omega=1$, the
flat factor is removed before taking the resultant.  Root orders and
theorems \ref{thm:jordan-regular}--\ref{thm:flat-band} distinguish
exceptional points from semisimple parity crossings without a numerical
tolerance.
\end{corollary}

\begin{proof}
The resultant vanishes exactly when the two displayed polynomials share
a root.  The preceding two theorems identify the Jordan block attached
to every such root and separately account for the flat factor.
\end{proof}

\section{Defect-localized modes}

\begin{definition}[Bulk spectral set]\label{def:bulk-set}
For the infinite translation-invariant chain, define
\[
 \Sigma_\omega=
 \left\{c-1-\omega\pm\sqrt{\omega^2-(1-c^2)}:
              c\in[-1,1]\right\}.
\]
For $\lambda\notin\Sigma_\omega$ with $a\ne0$, retain $a$ and $x$ from
\eqref{eq:a-def}--\eqref{eq:x-def}, and choose
\[
 s(\lambda)=\sqrt{x(\lambda)^2-1}
\]
as the standard branch on $x\in\C\setminus[-1,1]$ satisfying
$s/x\to1$ as $|x|\to\infty$.  Equivalently, it is the branch satisfying
$s(\lambda)\sim x(\lambda)$ as $|\lambda|\to\infty$.  Put $z=x-s$, so
$|z|<1$.
\end{definition}

\begin{theorem}[Infinite-line localization criterion]
\label{thm:localization}
Let $\delta\ne0$ and $\lambda\notin\Sigma_\omega$.  First suppose
$a\ne0$.  Then $\lambda$ is an exponentially localized eigenvalue of
the infinite chain with the same single tumbling defect if and only if
\begin{equation}\label{eq:localization-secular}
 1+2\delta g_\infty(\lambda)=0,
\end{equation}
where
\begin{equation}\label{eq:localization-green}
 g_\infty(\lambda)
 =\frac1{2a}\left(1+\frac{\lambda+1-x}{s}\right)
 =\frac{z(\lambda+1-z)}{a(1-z^2)}.
\end{equation}
The localization length is
\[
 \xi(\lambda)=\frac{-1}{\log|z|}.
\]
Equivalently, with $a=\lambda+1+\omega$, the pair $(a,z)$ satisfies
\begin{align}
 a^2+1-\omega^2-a(z+z^{-1})&=0,\label{eq:loc-dispersion}\\
 a(1-z^2)+2\delta z(a-\omega-z)&=0,
 \qquad 0<|z|<1.\label{eq:loc-match}
\end{align}
After elimination of $a$, every admissible $z$ is a root of
\begin{align}
0={}&2\delta z^3+(2\delta\omega+1-\omega^2)z^2\notag\\
&+(4\delta^2\omega+4\delta\omega^2-2\delta)z
 +(2\delta\omega+\omega^2-1),\label{eq:localization-cubic}
\end{align}
and, when the denominator is nonzero,
\begin{equation}\label{eq:localization-recovery}
 a=\frac{2\delta z(\omega+z)}{1-z^2+2\delta z},
 \qquad \lambda=a-1-\omega.
\end{equation}
A root of \eqref{eq:localization-cubic} is accepted only when
\eqref{eq:localization-secular}, or equivalently both
\eqref{eq:loc-dispersion} and \eqref{eq:loc-match}, is satisfied; this
removes roots introduced by elimination or by clearing denominators.

There is one additional compact case, at the point where $a=0$.  If
$\omega\ne1$, then
\begin{equation}\label{eq:compact-localization}
 \lambda=-1-\omega
 \quad\hbox{is localized if and only if}\quad
 \omega_0=\frac{\omega^2+1}{2\omega}.
\end{equation}
Its decay root is $z=0$, its localization length is $\xi=0$, and its
eigenvector is supported at the defect and its two neighboring sites.
Thus the admissible range in the preceding formulas is $0<|z|<1$;
$z=0$ is accepted precisely in \eqref{eq:compact-localization}.
\end{theorem}

\begin{proof}
The polarization entry of the bulk resolvent at the defect site is
\begin{align*}
 g_\infty(\lambda)
 &=\frac1{2\pi}\int_0^{2\pi}
   \frac{\lambda+1-\cos q}
        {a^2+1-\omega^2-2a\cos q}\,dq\\
 &=\frac1{2a}\left[1+(\lambda+1-x)
       \frac1{2\pi}\int_0^{2\pi}\frac{dq}{x-\cos q}\right]\\
 &=\frac1{2a}\left(1+\frac{\lambda+1-x}{s}\right).
\end{align*}
The rank-one eigenvalue equation is \eqref{eq:localization-secular}.  Since
$x=(z+z^{-1})/2$ and $s=(z^{-1}-z)/2$, elementary simplification gives
the second expression in \eqref{eq:localization-green}, and the Green
kernel decays as $z^{|j|}$.  Clearing the denominator in
\eqref{eq:localization-secular} gives
\eqref{eq:loc-match}; the bulk dispersion relation gives
\eqref{eq:loc-dispersion}.  Eliminating $a$ gives
\eqref{eq:localization-cubic}, while solving \eqref{eq:loc-match} for
$a$ gives \eqref{eq:localization-recovery}.

It remains to evaluate the point omitted by division by $a$.  At
$\lambda=-1-\omega$ the Fourier denominator is the nonzero constant
$1-\omega^2$, and hence
\[
 g_\infty(-1-\omega)
 =\frac1{2\pi}\int_0^{2\pi}
     \frac{-\omega-\cos q}{1-\omega^2}\,dq
 =\frac{\omega}{\omega^2-1}.
\]
The secular equation is therefore equivalent to
$\delta=(1-\omega^2)/(2\omega)$, which is exactly
\eqref{eq:compact-localization}.  In Fourier variables the associated
resolvent vector is proportional to
\[
 \binom{-\omega-e^{-\ii q}}{\omega+e^{\ii q}}.
\]
Its inverse Fourier transform is supported at the defect and its two
neighbors.  This proves the compact case and also shows directly that
it is the limiting root $z=0$.
\end{proof}

\begin{proposition}[Finite-ring convergence of a localized mode]
\label{prop:localized-convergence}
Suppose \eqref{eq:localization-secular} has a simple root
$\lambda_\infty$ with decay root
$z_\infty$, $0<|z_\infty|<1$.  Then, for all sufficiently large $n$, the
finite polynomial $Q_{n,-}$ has a unique root $\lambda_n$ near
$\lambda_\infty$, and
\[
 \lambda_n-\lambda_\infty=O(|z_\infty|^n).
\]
The corresponding right eigenvector has site mass bounded by
$C|z_\infty|^{2\operatorname{dist}(j,0)}$ after normalization.
In the compact case \eqref{eq:compact-localization}, the infinite and
finite eigenvalues agree exactly and the eigenvector is supported where
$\operatorname{dist}(j,0)\le1$.
\end{proposition}

\begin{proof}
The root-of-unity sum of the bulk Green kernel is the exact identity
\begin{equation}\label{eq:finite-green-z}
 G_n(\lambda)=\frac1{2a}\left[1+
  \frac{\lambda+1-x}{s}\frac{1+z^n}{1-z^n}\right].
\end{equation}
Indeed, Fourier expansion of $(x-\cos q)^{-1}$ gives
\[
 \frac1n\sum_{k=0}^{n-1}\frac1{x-\cos q_k}
 =\frac1s\frac{1+z^n}{1-z^n}.
\]
Both sides of \eqref{eq:finite-green-z} are meromorphic functions of
$\lambda$.  The identity therefore continues across the infinite-volume
cut to every point away from a finite-ring homogeneous pole; the two
boundary choices of $s$ give the same continued function.
Consequently
\begin{equation}\label{eq:finite-green-error}
 G_n-g_\infty=
 \frac{\lambda+1-x}{as}\frac{z^n}{1-z^n}.
\end{equation}
Choose a closed disk $B$ about $\lambda_\infty$ disjoint from
$\Sigma_\omega$ and containing no other zero of
$1+2\delta g_\infty$.  Shrinking it if necessary, choose
\[
 \rho=\max_{\lambda\in B}|z(\lambda)|<r<1.
\]
Equation
\eqref{eq:finite-green-error} and its $\lambda$-derivative are bounded
on $B$ by $C_Br^n$: for the derivative, the only additional factor is
$n\rho^{n-1}=O(r^n)$.  Rouch\'e's theorem on $\partial B$, followed by
the simple-zero estimate at $\lambda_\infty$, gives one zero
$\lambda_n$.  Evaluation of \eqref{eq:finite-green-error} at
$\lambda_\infty$, together with
$z(\lambda_n)=z_\infty+O(\lambda_n-\lambda_\infty)$, gives
$\lambda_n-\lambda_\infty=O(|z_\infty|^n)$.

The Fourier coefficients of the two bulk resolvent entries are linear
combinations of $z^{|j|}$ and $z^{|j\pm1|}$.  Their periodizations are
bounded by
\[
 \frac{r^{\operatorname{dist}(j,0)}
       +r^{n-\operatorname{dist}(j,0)}}{1-r^n}.
\]
The coefficients converge to their nonzero infinite-volume values, so
normalization changes this bound by a fixed factor.  Using
$z(\lambda_n)=z_\infty+O(|z_\infty|^n)$ gives the stated site-mass
estimate.

In the compact case, the Laurent-polynomial Fourier vector displayed in
the proof of Theorem \ref{thm:localization} is an exact finite-ring
eigenvector for every $n$ (with the two neighbors identified when
$n=2$), proving the last assertion.
\end{proof}

\section{Large-ring spectral gap}

\begin{definition}[Spectral gap]\label{def:gap}
For $\omega>0$ and $n\ge2$, irreducibility makes zero a simple
eigenvalue.  Define
\[
 \gamma_n(\omega,\omega_0)
 =-\max\{\Re\lambda:\lambda\in\spec L_n(\omega,\omega_0),
                            \ \lambda\ne0\}.
\]
Put $q=2\pi/n$ and
\begin{equation}\label{eq:hom-gap}
 \gamma_n^{\rm hom}(\omega)
 =1-\cos q+\omega-\sqrt{\omega^2-\sin^2q},
\end{equation}
where the square root is positive once $n$ is sufficiently large for
fixed $\omega>0$.
\end{definition}

\begin{lemma}[Protected first mode and affected pole]
\label{lem:gap-pole}
For every $n$, the protected sector contains
$\lambda_1^+=-\gamma_n^{\rm hom}$.  For fixed $\omega>0$ and
$\omega_0\ge0$, the affected root born from the same homogeneous pole
is
\begin{equation}\label{eq:affected-gap-shift}
 \lambda_{1,-}=\lambda_1^+
 +\frac{\delta(1+\omega)}{\omega^2(1+\omega_0)}
       \frac{q^2}{n}+O(|\delta|n^{-4}).
\end{equation}
\end{lemma}

\begin{proof}
The first assertion is Theorem \ref{thm:parity-factor}.  Near
$\lambda_1^+$, the two terms $k=1,n-1$ in the full Fourier Green
function have residue
\[
 R_1=\frac{\lambda_1^++1-c_1}{\lambda_1^+-\lambda_1^-}
    =-\frac{\omega-r_1}{2r_1},
 \qquad r_1=\sqrt{\omega^2-\sin^2q}.
\]
Thus
\[
 G_n(\lambda)=\frac{2R_1}{n(\lambda-\lambda_1^+)}
                  +G_n^{\rm reg}(\lambda).
\]
At $\lambda=0$, every nonzero-wave summand is exactly
$1/[2(1+\omega)]$.  Splitting the sum into
$|q_k|\le n^{-1/2}$ and its complement and using
$\lambda_1^+=O(n^{-2})$ gives
\[
 G_n^{\rm reg}(\lambda_1^+)
 =\frac1{2(1+\omega)}+O(n^{-1}),
 \qquad
 (G_n^{\rm reg})'(\lambda)=O(n)
\]
for $|\lambda-\lambda_1^+|=O(n^{-3})$.  Solving
$1+2\delta G_n=0$ gives
\[
 \lambda_{1,-}-\lambda_1^+
 =\frac{2\delta(\omega-r_1)}{nr_1}
    \frac{1}{1+\delta/(1+\omega)}+O(|\delta|n^{-4}),
\]
which expands to \eqref{eq:affected-gap-shift}.  The proportional
remainder follows because the secular equation is identical to the
homogeneous one at $\delta=0$; tracking the estimates for the regular
part above leaves a factor $\delta$ in every omitted term.
\end{proof}

\begin{lemma}[Right-edge root separation]\label{lem:right-edge-separation}
Fix $\omega>0$ and $\omega_0\ge0$, and put
$D=(1+\omega)/(2\omega)$.  There is $N$ such that, for
$n\ge N$, every nonzero eigenvalue other than the protected and affected
roots born from $q=\pm2\pi/n$ satisfies
\begin{equation}\label{eq:right-edge-separation}
 \Re\lambda\le \Re\lambda_1^+-\pi^2D n^{-2}.
\end{equation}
The constants can be chosen uniformly when $(\omega,\omega_0)$ ranges
over a compact set with $\omega$ bounded away from zero.
\end{lemma}

\begin{proof}
Put
\[
 \Omega_n=\left\{\lambda:\Re\lambda\ge
 \Re\lambda_1^+-\pi^2D n^{-2}\right\}.
\]
The protected roots are explicit.  Their slow branch satisfies, uniformly
for small real $q$,
\begin{equation}\label{eq:slow-dispersion-bound}
 \lambda^+(q)=-Dq^2+O(q^4).
\end{equation}
Choose $q_*>0$ so that the remainder is at most $Dq^2/4$ for
$|q|\le q_*$.  The modes with $2\pi/n<|q|\le q_*$ are then to the left
of $\Omega_n$, while continuity puts both branches with
$q_*\le|q|\le\pi$ a fixed distance into the left half-plane.  The fast
small-$q$ branch is also separated from zero by a fixed amount.  Thus,
apart from zero, only the protected root born from $q=\pm2\pi/n$ lies
in $\Omega_n$ once $n$ is large.  We prove the same exclusion for the
affected sector.

If $\delta=0$, that sector is homogeneous and the conclusion follows from
\eqref{eq:slow-dispersion-bound}.  Assume $\delta\ne0$ and set
 \begin{equation}\label{eq:threshold-secular-margin}
 f_n(\lambda)=1+2\delta G_n(\lambda),\qquad
 g_0=\frac1{2(1+\omega)},\qquad
 M_0=1+2\delta g_0=\frac{1+\omega_0}{1+\omega}>0.
 \end{equation}
In a neighborhood of zero, every affected root is a zero of $f_n$.
Indeed, the residue at every slow homogeneous pole is nonzero, so
multiplication by $H_{n,-}^{(0)}$ cancels the pole of $f_n$ but does not
leave a zero there.

We first exclude affected zeros that do not tend to zero.  Gershgorin's
theorem places every finite-ring eigenvalue in a fixed disk and in the
closed left half-plane.  If the desired exclusion failed, there would
be $\epsilon>0$ and a sequence of affected zeros in $\Omega_n$ with
$|\lambda|\ge\epsilon$; a subsequence would converge to
$\lambda_*\ne0$ with $\Re\lambda_*=0$.  The bulk spectral set meets the
imaginary axis only at zero, so \eqref{eq:finite-green-error} gives
$G_n\to g_\infty$ locally uniformly at $\lambda_*$.  The limiting
secular equation would produce an $\ell^2$ eigenvector $u$ of the
infinite defective generator with eigenvalue $\lambda_*$.  This is
impossible because
\begin{equation}\label{eq:infinite-dissipation}
 -\Re\langle u,L_\infty u\rangle
 =\frac12\sum_{j,\sigma}|u_{j+1}^\sigma-u_j^\sigma|^2
   +\sum_j\omega_j|u_j^+-u_j^-|^2.
\end{equation}
Equality in \eqref{eq:infinite-dissipation} forces an $\ell^2$ vector to
vanish.  The subsequence contradiction therefore supplies
$\epsilon>0$ such that no affected zero in $\Omega_n$ can satisfy
$|\lambda|\ge\epsilon$.

It remains to exhaust the zeros tending to zero.  Put $\lambda=\zeta/n^2$
and, using meromorphic continuation of \eqref{eq:finite-green-z}, choose
either local branch of
\[
 w_n(\zeta)=n\operatorname{arccosh}x(\zeta/n^2),
 \qquad z(\zeta/n^2)^n=e^{-w_n(\zeta)}.
\]
The quantities $w_n^2$ and $w_n\coth(w_n/2)$ below are even in $w_n$,
so the choice does not affect the argument, including on the continued
bulk cut.
The definitions of $a$ and $x$ give
\begin{equation}\label{eq:fixed-threshold-expansions}
 x-1=\frac{\omega}{1+\omega}\lambda+O(\lambda^2),
 \qquad
 \lambda+1-x=\frac{\lambda}{1+\omega}+O(\lambda^2).
\end{equation}
Since $s=\sinh(w_n/n)$, equations
\eqref{eq:finite-green-z} and \eqref{eq:fixed-threshold-expansions} give the
uniform expansion on bounded $\zeta$-sets away from its poles
\begin{align}
 w_n(\zeta)^2&=\frac{\zeta}{D}+O(n^{-2}),
 \label{eq:fixed-scaled-w}\\
 G_n(\zeta/n^2)
 &=g_0+\frac{w_n(\zeta)}{4\omega(1+\omega)n}
       \coth\!\left(\frac{w_n(\zeta)}2\right)+O(n^{-2}).
 \label{eq:fixed-scaled-green}
\end{align}
The error in \eqref{eq:fixed-scaled-green} remains $O(n^{-2})$ when the
distance from a nonzero pole of the hyperbolic cotangent is at least
$A/n$, for fixed $A>0$.  The poles in the limiting $\zeta$-plane are
$-4\pi^2Dk^2$.  The scaled half-plane $n^2\Omega_n$ contains only the
first one, $-4\pi^2D$.  Near that pole the second term in
\eqref{eq:fixed-scaled-green} is bounded by $C/A$ outside
$|\lambda-\lambda_1^+|<A n^{-3}$.  Choose $A$ so that
$2|\delta|C/A<M_0/4$.  After increasing $n$,
\eqref{eq:fixed-scaled-green} gives
\begin{equation}\label{eq:fixed-exclusion-margin}
 |f_n(\lambda)|\ge \frac{M_0}{2}
\end{equation}
there, outside a disk
$|\lambda-\lambda_1^+|<A n^{-3}$, provided $A$ and then $n$ are large
enough.  If instead $\lambda\in\Omega_n$ tends to zero while
$n^2|\lambda|\to\infty$, then its imaginary part dominates its negative
real part.  Since $x-1=\omega\lambda/(1+\omega)+O(\lambda^2)$, the
standard square-root expansion of $\operatorname{arccosh}$ gives
$\Re\operatorname{arccosh}x(\lambda)\ge c\sqrt{|\lambda|}$.
Consequently $z(\lambda)^n\to0$, and
\eqref{eq:finite-green-z} again gives $G_n(\lambda)=g_0+o(1)$ and
\eqref{eq:fixed-exclusion-margin}.  These two cases exhaust the region
near zero.

Finally, on the circle $|\lambda-\lambda_1^+|=A n^{-3}$ write
\[
 G_n(\lambda)=\frac{2R_1}{n(\lambda-\lambda_1^+)}
                 +G_n^{\rm reg}(\lambda).
\]
Equations \eqref{eq:fixed-scaled-green} and
\eqref{eq:threshold-secular-margin} give
$1+2\delta G_n^{\rm reg}=M_0+O(n^{-1})$.  After multiplication by
$\lambda-\lambda_1^+$, the secular equation becomes
\begin{equation}\label{eq:fixed-first-pole-rouche}
 (\lambda-\lambda_1^+)f_n(\lambda)
 =M_0(\lambda-\lambda_1^+)+\frac{4\delta R_1}{n}
   +O(n^{-1}|\lambda-\lambda_1^+|).
\end{equation}
Here $R_1=O(n^{-2})$.  Enlarge $A$, if necessary, so that the zero of
the first two terms lies strictly inside the circle.  On the circle the
last term is $O(n^{-4})$, whereas the first two terms have modulus
bounded below by $c_A n^{-3}$.  Rouch\'e's theorem therefore counts
exactly one zero in the disk.  It is the affected root of Lemma
\ref{lem:gap-pole}.  Thus no other affected root occurs in $\Omega_n$,
which proves the assertion.

All estimates above are uniform on the stated compact parameter sets:
$D$, $\delta$, and the generator norms are bounded there, while $M_0$
has a positive minimum.
\end{proof}

\begin{theorem}[Fixed-rate gap asymptotics]\label{thm:fixed-gap}
Fix $\omega>0$ and $\omega_0\ge0$.
\begin{enumerate}
\item If $\omega_0\le\omega$, then for every sufficiently large $n$,
  \begin{equation}\label{eq:slow-defect-gap}
    \gamma_n(\omega,\omega_0)=\gamma_n^{\rm hom}(\omega).
  \end{equation}
\item If $\omega_0>\omega$, then
  \begin{equation}\label{eq:fast-defect-gap}
 \gamma_n(\omega,\omega_0)
  =\gamma_n^{\rm hom}(\omega)
   -\frac{4\pi^2(\omega_0-\omega)(1+\omega)}
          {\omega^2(1+\omega_0)}\,n^{-3}
    +O(|\omega_0-\omega|n^{-4}).
  \end{equation}
\item In both cases,
  \begin{align}
  \gamma_n^{\rm hom}(\omega)
   ={}&\frac{2\pi^2(1+\omega)}{\omega}\,n^{-2}\notag\\
    &-16\pi^4\left(\frac1{24}+\frac1{6\omega}
                    -\frac1{8\omega^3}\right)n^{-4}
      +O(n^{-6}).\label{eq:hom-gap-expansion}
  \end{align}
\end{enumerate}
The errors are uniform when $(\omega,\omega_0)$ ranges over a compact
set with $\omega$ bounded away from zero.  The root selection is uniform
through $\omega_0=\omega$, where the two parity roots coincide.
\end{theorem}

\begin{proof}
Taylor expansion of \eqref{eq:hom-gap} gives
\eqref{eq:hom-gap-expansion}.  Lemma \ref{lem:gap-pole} shows
that, uniformly across $\delta=0$, the affected first root moves toward
zero precisely when $\delta>0$.  Lemma
\ref{lem:right-edge-separation} shows that no other
extended, band-edge, or localized root can compete with this first
pair.  Therefore the protected root wins
when $\delta\le0$, proving \eqref{eq:slow-defect-gap}, and the affected
root wins when
$\delta>0$.  Substitution of $q^2=4\pi^2n^{-2}$ in
\eqref{eq:affected-gap-shift} then gives \eqref{eq:fast-defect-gap}.
\end{proof}

\begin{corollary}[Fixed ratio, vanishing ratio, and a slow-site scaling]
\label{cor:requested-regimes}
Let the bulk rate $\omega>0$ be fixed.
\begin{enumerate}
\item If $\omega_0/\omega\to\rho\in(0,\infty)$, equations
  \eqref{eq:slow-defect-gap}--\eqref{eq:hom-gap-expansion} apply.  For
  $\rho<1$ the protected branch is selected; for
  $\rho>1$ the affected branch and its $n^{-3}$ term are selected.  At
  the boundary $\rho=1$, the protected root is selected whenever
  $\omega_0\le\omega$ and the affected root whenever
  $\omega_0>\omega$, with the expansion uniform across equality.
\item If $\omega_0/\omega\to0$, then \eqref{eq:slow-defect-gap} holds
  for every sufficiently large $n$.
\item In particular, if $n\omega_0\to\beta$ with $\beta\ge0$, then
  \eqref{eq:slow-defect-gap} still holds.  A single slow tumbling site
  therefore has no
  separate $\omega_0\asymp n^{-1}$ gap crossover when the bulk rate is
  fixed.
\end{enumerate}
\end{corollary}

\begin{proof}
In items 2 and 3 one has $\omega_0<\omega$ for every sufficiently large
$n$, so Theorem \ref{thm:fixed-gap} applies.  Item 1 follows by the same
sign test and compact-uniform expansion.
\end{proof}

\begin{lemma}[Critical right-edge exhaustion]
\label{lem:critical-exhaustion}
Set $\omega=\alpha/n$ and $\omega_0=\beta/n$, where
$\alpha>0$ and $\beta\ge0$ are fixed, and put $k=2\pi$.  For all
sufficiently large $n$, every nonzero eigenvalue with maximal real part
belongs to the protected or affected cluster born from
$q=\pm k/n$.  More precisely, if $\Lambda_{1,n}$ is the largest real
part within that cluster, then every eigenvalue outside the cluster
satisfies
\begin{equation}\label{eq:critical-exhaustion-separation}
 \Re\lambda\le\Lambda_{1,n}-
 \begin{cases}
  c n^{-2},&\alpha\le k,\\
  c n^{-1},&\alpha>k,
 \end{cases}
\end{equation}
for a constant $c=c(\alpha,\beta)>0$.
\end{lemma}

\begin{proof}
The protected roots are the homogeneous roots, so only the affected
sector requires an exhaustion argument.  If $\beta=\alpha$, the two
parity restrictions are homogeneous, the result follows directly from
their explicit dispersion, and \eqref{eq:critical-uniform-simple-shift}
below holds with zero shift.  Assume henceforth that
$b:=\beta-\alpha\ne0$.  Write
\[
 p_{\ell,n}^{\pm}=\cos(\ell k/n)-1-\frac\alpha n
 \pm\sqrt{\frac{\alpha^2}{n^2}-\sin^2(\ell k/n)}
\]
for its homogeneous poles.  Away from a double pole, the paired
$q=\pm\ell k/n$ contribution to $G_n$ has the partial-fraction form
\begin{equation}\label{eq:critical-pole-fraction}
 \frac{2R_{\ell,n}^\pm}
      {n(\lambda-p_{\ell,n}^\pm)},
 \qquad
 R_{\ell,n}^\pm=
 \frac{p_{\ell,n}^\pm+1-\cos(\ell k/n)}
      {p_{\ell,n}^\pm-p_{\ell,n}^\mp}.
\end{equation}

We first prove an off-pole estimate that exhausts the affected roots.
Gershgorin's theorem confines all eigenvalues to a fixed disk, so below
we restrict the right-edge windows to that disk.  If
$\alpha\le k$, choose $B>0$ large enough to contain the protected first
root and, when $\alpha<k$, its affected $O(n^{-2})$ displacement scale.
We take $B$ midway between two consecutive fixed-index coefficients
$\ell^2k^2/2$ and consider
\[
 \Omega_{B,n}=\left\{\Re\lambda\ge
             -\frac\alpha n-\frac B{n^2}\right\};
\]
if $\alpha>k$, choose $d$ strictly between
\[
 d_1=\alpha-\sqrt{\alpha^2-k^2},\qquad
 d_2=\alpha-\Re\sqrt{\alpha^2-4k^2},
\]
and put $\Omega_{d,n}=\{\Re\lambda\ge-d/n\}$.  Only finitely many
fixed-index pole centers lie in $\Omega_{B,n}$: indeed, their indices
satisfy $1-\cos(\ell k/n)\le Bn^{-2}$.  In $\Omega_{d,n}$ only the slow
first pole can occur for all sufficiently large $n$.  Remove the disks
of radius $A n^{-2}$ about these simple poles.  When $\alpha=k$, replace
the two first-pole disks by the single disk of radius $A n^{-3/2}$ about
$(p_{1,n}^++p_{1,n}^-)/2$.

On the remaining set the exact Fourier sum satisfies
\begin{equation}\label{eq:critical-green-exclusion-bound}
 |G_n(\lambda)|\le C_0(B)+C_1\log n
             +C_2(B)\frac nA+C_3\frac n{\sqrt B}
 \qquad(\alpha\le k),
\end{equation}
where $C_1,C_3$ are independent of $A,B,n$ once $B$ exceeds a fixed
threshold.  For $\alpha>k$ the same bound holds without the
$C_3n/\sqrt B$ term and with the remaining constants depending on $d$.
Here is the estimate.  First combine the
two roots belonging to each quadratic $D_\ell$; this cancellation is
essential.  For $|q_\ell|\le q_0$, the poles outside the finite set have
horizontal distance at least
$c\ell^2n^{-2}-Bn^{-2}$ from $\Omega_{B,n}$.  After scaling imaginary
parts by $n$, choose the pole whose scaled imaginary part is closest to
$n\Im\lambda$ and call its index $\ell_0$.  Its contribution is at most
$Cn/\sqrt B$; the midpoint choice of $B$ gives the required coefficient
gap at the first exterior index.  The $j$th neighboring index contributes
at most $C/(1+j)$.
This gives the last and logarithmic terms in
\eqref{eq:critical-green-exclusion-bound}.  Each removed simple-pole
term is at most $Cn/A$ off its disk.  At $\alpha=k$, the exact grouped
first contribution, with
$\Delta=\lambda-(p_{1,n}^++p_{1,n}^-)/2$, is
\[
 \frac{2(\Delta-k/n)}
      {n\{\Delta^2-(p_{1,n}^+-p_{1,n}^-)^2/4\}},
\]
and is at most $Cn/A^2$ outside the $A n^{-3/2}$ disk.  Finally, the
terms with $|q_\ell|\ge q_0$ have denominators bounded away from zero
in these right-edge windows and form a bounded Riemann sum.  These four
estimates prove \eqref{eq:critical-green-exclusion-bound}.

Since $2\delta=2b/n$, first enlarge $B$ so that
$2|b|C_3/\sqrt B<1/8$, then enlarge $A$ so that
$2|b|C_2(B)/A<1/8$, and finally increase $n$ so that the constant and
logarithmic terms contribute less than $1/4$.  In the case $\alpha>k$,
choose $A$ first and then $n$.  Thus
\eqref{eq:critical-green-exclusion-bound} implies
\begin{equation}\label{eq:critical-secular-exclusion}
 |2\delta G_n(\lambda)|<\frac12.
\end{equation}
Thus every zero of $1+2\delta G_n$ in a right-edge window lies in one of
the displayed pole neighborhoods.

We also need the root count inside a simple-pole disk.  There
\[
 G_n(\lambda)=\frac{2R_{\ell,n}^\pm}
  {n(\lambda-p_{\ell,n}^\pm)}+G_{\ell,n}^{\rm reg}(\lambda),
 \qquad G_{\ell,n}^{\rm reg}=O(1).
\]
The last bound follows directly from \eqref{eq:Gfinite}: for the
remaining small wave numbers the combined quadratic summand is bounded
by $C/(1+|j^2-\ell^2|)$ up to $j\asymp\sqrt n$, while each of the
remaining $O(n)$ summands is $O(n^{-1})$.  Multiplication of the secular
equation by $\lambda-p_{\ell,n}^\pm$ gives
\begin{equation}\label{eq:critical-simple-rouche}
 (\lambda-p_{\ell,n}^\pm)(1+2\delta G_n)
 =(\lambda-p_{\ell,n}^\pm)
   \left(1+\frac{2b}{n}G_{\ell,n}^{\rm reg}\right)
   +\frac{4bR_{\ell,n}^\pm}{n^2}.
\end{equation}
Compare the right-hand side with
\[
 h_0(\lambda)=\lambda-p_{\ell,n}^\pm
                  +4bR_{\ell,n}^\pm n^{-2}.
\]
Enlarge $A$ once more so that the zero of $h_0$ is strictly inside the
disk.  On its boundary, the difference from $h_0$ is $O(n^{-3})$,
whereas $|h_0|\ge c_A n^{-2}$.  Rouch\'e's theorem gives exactly one
affected zero and
\begin{equation}\label{eq:critical-uniform-simple-shift}
 \widehat p_{\ell,n}^\pm
 =p_{\ell,n}^\pm
  -4bR_{\ell,n}^\pm n^{-2}+O(n^{-3}).
\end{equation}
Together with \eqref{eq:critical-secular-exclusion}, this accounts for
every affected root in either right-edge window.

Suppose first that $\alpha<k$.  For every fixed $\ell\ge1$, the two
poles are nonreal conjugates and
\begin{equation}\label{eq:critical-residue-real-half}
 \Re R_{\ell,n}^\pm=\frac12,
 \qquad
 \Re p_{\ell,n}^\pm
 =-\frac\alpha n-\frac{\ell^2k^2}{2n^2}+O(n^{-4}).
\end{equation}
Equations \eqref{eq:critical-uniform-simple-shift}--
\eqref{eq:critical-residue-real-half} show that the same parity is
selected for every fixed $\ell$, and that the $\ell=2$ cluster is to the
left of the first by $3k^2/(2n^2)+O(n^{-3})$.  The choice of $B$ puts
every possible competitor either among these finitely many clusters or
to their left and separates all roots outside the window from the first
cluster by a fixed multiple of $n^{-2}$.  The endpoint tumble root has
real part $-2\alpha/n+O(n^{-2})$.  This proves
\eqref{eq:critical-exhaustion-separation} for $\alpha<k$.

At $\alpha=k$, the first protected cluster contains
\begin{equation}\label{eq:critical-protected-reference}
 -\frac{k}{n}
 +\left(-\frac{k^2}{2}+\frac{k^2}{\sqrt3}\right)n^{-2}
 +O(n^{-3}).
\end{equation}
For every fixed $\ell\ge2$, equations
\eqref{eq:critical-uniform-simple-shift}--
\eqref{eq:critical-residue-real-half} apply with $\alpha=k$.  If
$\beta\le k$, its largest possible real part is below
\eqref{eq:critical-protected-reference} by
\[
 \left[\frac{(\ell^2-1)k^2}{2}
       +\frac{k^2}{\sqrt3}-2(k-\beta)\right]n^{-2}
 +O(n^{-3}),
\]
which is positive because $k>2\sqrt3$.  If $\beta>k$, the affected
higher modes move to the left and the protected comparison is even
stronger.  Estimate \eqref{eq:critical-secular-exclusion} accounts for
all remaining affected roots, proving the assertion at $\alpha=k$.

Finally, if $\alpha>k$, the first slow pole has decay coefficient
$d_1=\alpha-\sqrt{\alpha^2-k^2}$.  By the choice of $d$, the only
protected pole in $\Omega_{d,n}$ is the first slow pole, and
\eqref{eq:critical-secular-exclusion}--
\eqref{eq:critical-uniform-simple-shift} show that its affected partner
is the only affected root there.  Both have real part
$-d_1/n+O(n^{-2})$, whereas every root outside the window has real part
at most $-d/n$.  This gives a separation of order $n^{-1}$ and completes
the proof.
\end{proof}

\begin{theorem}[Bulk-critical double scaling]\label{thm:critical-gap}
Set
\[
 \omega=\frac{\alpha}{n},\qquad
 \omega_0=\frac{\beta}{n},\qquad
 k=2\pi,
\]
where $\alpha>0$ and $\beta\ge0$ are fixed.  Away from $\alpha=k$,
\begin{align}
\alpha<k:\quad
 \gamma_n
 &=\frac{\alpha}{n}
   +\left[\frac{k^2}{2}-2(\alpha-\beta)_+\right]n^{-2}
   +O(n^{-3}),\label{eq:critical-below}\\
\alpha>k:\quad
 \gamma_n
 &=\frac{\alpha-\sqrt{\alpha^2-k^2}}{n}\notag\\
 &\quad+\left[\frac{k^2}{2}
 -2(\beta-\alpha)_+
   \frac{\alpha-\sqrt{\alpha^2-k^2}}
        {\sqrt{\alpha^2-k^2}}\right]n^{-2}
 +O(n^{-3}).\label{eq:critical-above}
\end{align}
At $\alpha=k$,
\begin{align}
 \beta\le k:\quad
 \gamma_n&=\frac{k}{n}
 +k^2\left(\frac12-\frac1{\sqrt3}\right)n^{-2}
 +O(n^{-3}),\label{eq:critical-at-slow}\\
 \beta>k:\quad
 \gamma_n&=\frac{k}{n}
 -2\sqrt{k(\beta-k)}\,n^{-3/2}+O(n^{-2}).
 \label{eq:critical-at-fast}
\end{align}
Thus the first homogeneous exceptional scaling is a square-root boundary
layer for a faster defect, while the side selecting the affected parity
mode reverses as $\alpha$ passes through $2\pi$.
\end{theorem}

\begin{proof}
For $\alpha<k$, the protected first pair has
\[
 \Re\lambda_1^+=-\frac{\alpha}{n}-\frac{k^2}{2n^2}
 +O(n^{-3}).
\]
The pole residue tends to
$R_1=\tfrac12+\ii\alpha/(2\sqrt{k^2-\alpha^2})$, and
$\delta=(\beta-\alpha)/n$.  Hence the affected real part shifts by
$-2(\beta-\alpha)n^{-2}+O(n^{-3})$ by
\eqref{eq:critical-uniform-simple-shift}.  Taking the larger real part
gives \eqref{eq:critical-below}.

For $\alpha>k$, put $s=\sqrt{\alpha^2-k^2}$.  The protected rate is
$(\alpha-s)n^{-1}+k^2(2n^2)^{-1}+O(n^{-3})$, while
$R_1\to-(\alpha-s)/(2s)$.  The same pole equation shifts the affected
eigenvalue by
$2(\beta-\alpha)(\alpha-s)/(sn^2)+O(n^{-3})$.  Selection of the larger
real part gives \eqref{eq:critical-above}.  Lemma
\ref{lem:critical-exhaustion} shows in both cases that no root outside
the first cluster can compete with the selected root.

At $\alpha=k$, the exact lattice splitting of the homogeneous first
block is
\[
 \sqrt{\omega^2-\sin^2(k/n)}
   =\frac{k^2}{\sqrt3\,n^2}+O(n^{-4}),
\]
which proves \eqref{eq:critical-at-slow} when the protected root wins.
Near the center of this
block, put
\[
 \lambda_c=\cos(k/n)-1-k/n,\qquad
 r_n=\sqrt{(k/n)^2-\sin^2(k/n)},\qquad
 \Delta=\lambda-\lambda_c.
\]
The two $q=\pm k/n$ terms can be retained exactly, while the remaining
terms are uniformly bounded for $|\Delta|\asymp n^{-3/2}$:
\begin{equation}\label{eq:critical-green-two-terms}
 G_n(\lambda)
 =\frac{2(\Delta-k/n)}{n(\Delta^2-r_n^2)}+O(1)
 =-\frac{2k}{n^2\Delta^2}+\frac{2}{n\Delta}+O(1).
\end{equation}
For completeness, the bounded remainder follows by pairing wave numbers
$\ell$ and $n-\ell$: for each fixed $\ell\ne1$ its magnitude is bounded
by a constant times $(\ell^2-1)^{-1}$, the sum of these bounds
converges, and the terms with $|q_\ell|$ bounded away from zero form a
bounded Riemann sum.
For $\beta>k$, put $b=\beta-k$.  Substitution of
$\Delta=n^{-3/2}y$ in $1+2bG_n/n=0$ gives
$y^2=4kb+O(n^{-1/2})$.  Rouch\'e's theorem on disks of radius
$C n^{-2}$ about the two approximations, with fixed sufficiently large
$C$, gives two affected roots
\[
 \Delta=\mathord\pm2\sqrt{kb}\,n^{-3/2}+O(n^{-2}).
\]
The positive displacement gives \eqref{eq:critical-at-fast}.
For $\beta<k$, put $h=k-\beta>0$ and substitute
\[
 \Delta=n^{-3/2}\bigl(y+xn^{-1/2}\bigr)
\]
in $1-2hG_n/n=0$.  The constant and $n^{-1/2}$ coefficients in
\eqref{eq:critical-green-two-terms} give, respectively,
\[
 y^2=-4kh,\qquad x=2h.
\]
Rouch\'e's theorem on two disks of radius $C n^{-5/2}$ about the
resulting approximations, with fixed sufficiently large $C$, gives the
affected pair
\begin{equation}\label{eq:critical-slow-affected-pair}
 \Delta=\mathord\pm2\ii\sqrt{kh}\,n^{-3/2}
        +2h n^{-2}+O(n^{-5/2}).
\end{equation}
Its decay rate exceeds the protected decay rate by
\[
 \left(\frac{k^2}{\sqrt3}-2h\right)n^{-2}+O(n^{-5/2}).
\]
This coefficient is positive for $0<h\le k$ because
$k=2\pi>2\sqrt3$.  The protected root therefore wins for every
$\beta<k$; at $\beta=k$ the two parity roots coincide.  This proves
\eqref{eq:critical-at-slow}; the case $\beta>k$ was obtained from the
positive displacement above.  Lemma
\ref{lem:critical-exhaustion} supplies the required exclusion of every
higher, endpoint, and off-pole affected root in both cases.
\end{proof}

\section{Telegraph comparison}

\begin{definition}[Two continuum approximations]\label{def:telegraph}
For a homogeneous ring of circumference $n$, the standard telegraph
system is
\begin{equation}\label{eq:telegraph-system}
 \partial_t\rho=-\partial_xm,
 \qquad
 \partial_tm=-\partial_x\rho-2\omega m.
\end{equation}
Its slow Fourier eigenvalue at wave number $q$ is
\begin{equation}\label{eq:telegraph-eigenvalue}
 \lambda_{\rm tel}^+(q)=-\omega+\sqrt{\omega^2-q^2}.
\end{equation}
The lattice-diffusion-corrected system is
\begin{equation}\label{eq:corrected-telegraph-system}
 \partial_t\rho=-\partial_xm+\tfrac12\partial_x^2\rho,
 \qquad
 \partial_tm=-\partial_x\rho+\tfrac12\partial_x^2m-2\omega m.
\end{equation}
Its slow eigenvalue is
\begin{equation}\label{eq:corrected-telegraph-eigenvalue}
 \lambda_{\rm td}^+(q)=-\frac{q^2}{2}-\omega
                         +\sqrt{\omega^2-q^2}.
\end{equation}
\end{definition}

\begin{theorem}[Exact relaxation rate versus telegraph rates]
\label{thm:telegraph-comparison}
Fix $\omega>0$, let $q=2\pi/n$, and define
$\gamma_{\rm tel}=-\lambda_{\rm tel}^+(q)$ and
$\gamma_{\rm td}=-\lambda_{\rm td}^+(q)$.  Then
\begin{align}
 \gamma_n^{\rm hom}-\gamma_{\rm tel}
 &=\frac{q^2}{2}
   -\left(\frac1{24}+\frac1{6\omega}\right)q^4+O(q^6),
   \label{eq:telegraph-error}\\
 \gamma_n^{\rm hom}-\gamma_{\rm td}
 &=-\left(\frac1{24}+\frac1{6\omega}\right)q^4+O(q^6).
   \label{eq:corrected-telegraph-error}
\end{align}
Consequently
\begin{equation}\label{eq:telegraph-ratios}
 \frac{\gamma_{\rm tel}}{\gamma_n^{\rm hom}}
 \longrightarrow\frac1{1+\omega},
 \qquad
 \frac{\gamma_{\rm td}}{\gamma_n^{\rm hom}}\longrightarrow1.
\end{equation}
In the bulk-critical scaling $\omega=\alpha/n$, the standard telegraph
rate does recover the leading term:
\begin{equation}\label{eq:telegraph-critical}
 \gamma_{\rm tel}
 =\frac{\alpha-\Re\sqrt{\alpha^2-(2\pi)^2}}{n},
\end{equation}
which agrees with the leading terms of
\eqref{eq:critical-below}--\eqref{eq:critical-at-fast}.
\end{theorem}

\begin{proof}
The exact lattice dispersion replaces $-q^2/2$ by $\cos q-1$ and $q$
inside the square root by $\sin q$.  Expanding
$\cos q$, $\sin^2q$, and the square root gives
\eqref{eq:telegraph-error}--\eqref{eq:corrected-telegraph-error}.
Dividing by the leading exact rate $(1+\omega)q^2/(2\omega)$ gives
\eqref{eq:telegraph-ratios}.  Substitution of $\omega=\alpha/n$ and
$q=2\pi/n$ into \eqref{eq:telegraph-eigenvalue} gives
\eqref{eq:telegraph-critical}.
\end{proof}

\begin{proposition}[Point-defect telegraph secular equation]
\label{prop:telegraph-defect}
On the circle cut at zero, define the point-defect telegraph operator by
\[
 {\cal A}_\delta(\rho,m)=(-m',-\rho'-2\omega m)
 \quad\hbox{on }(0,n),
\]
with domain consisting of $m,\rho\in H^1(0,n)$ having traces that
satisfy
\begin{equation}\label{eq:telegraph-jump-domain}
 m(0+)=m(n-),\qquad
 \rho(0+)-\rho(n-)=-2\delta m(0).
\end{equation}
This is the distributional realization obtained by adding
$-2\delta\,\delta_0(x)m$ to the second equation in
\eqref{eq:telegraph-system}.  Outside the homogeneous telegraph
spectrum, put
\[
 \kappa=\sqrt{\lambda(\lambda+2\omega)}
\]
and interpret
$\kappa^{-1}\coth(n\kappa/2)$ as the even meromorphic function of
$\kappa$ obtained by analytic continuation; no sign choice for the
square root is required.  The affected-parity eigenvalues satisfy
\begin{equation}\label{eq:telegraph-defect}
 1+\delta\frac{\lambda}{\kappa}
            \coth\left(\frac{n\kappa}{2}\right)=0.
\end{equation}
The reflection-even telegraph modes are unchanged.  Formula
\eqref{eq:telegraph-defect} is the
continuum counterpart of \eqref{eq:resolvent-form}; the exact lattice
formula retains $1-\cos q$ and $\sin q$, and
\eqref{eq:telegraph-error}--\eqref{eq:corrected-telegraph-error} quantify
the resulting relaxation-rate error.
\end{proposition}

\begin{proof}
The jump in \eqref{eq:telegraph-jump-domain} cancels the delta
distribution in $-\rho'-2\delta\delta_0m(0)$, while continuity of $m$
prevents a delta distribution in $-m'$.  Thus the stated domain gives
the point interaction on $L^2((0,n);\mathbb C^2)$.  It is densely
defined because it contains $C_c^\infty((0,n);\mathbb C^2)$, and it is
closed:
the graph norm controls $\|m'\|_2$ directly and controls
$\|\rho'\|_2$ by
$\|\rho'\|_2\le\|({\cal A}_\delta(\rho,m))_2\|_2+
2\omega\|m\|_2$; the trace theorem then preserves
\eqref{eq:telegraph-jump-domain} under graph-norm limits.

Let $\lambda$ lie outside the homogeneous telegraph spectrum.  The
eigenvalue equations away from zero are
\[
 -m'=\lambda\rho,\qquad -\rho'-2\omega m=\lambda m.
\]
Thus $\lambda\ne0$ and elimination of $\rho$ gives
$m''=\kappa^2m$.  Write
\[
 m(x)=A\cosh\!\left(\kappa(x-n/2)\right)
      +B\sinh\!\left(\kappa(x-n/2)\right).
\]
The equality $m(0)=m(n)$ gives $B=0$, since
$\sinh(n\kappa/2)=0$ would put $\lambda$ in the homogeneous spectrum.
Using $\rho=-m'/\lambda$, the second jump condition in
\eqref{eq:telegraph-jump-domain} is equivalent to
\begin{equation}\label{eq:telegraph-derivative-jump}
 m'(0)-m'(n)=2\delta\lambda m(0).
\end{equation}
Substitution of the even hyperbolic-cosine solution into
\eqref{eq:telegraph-derivative-jump} gives
\[
 -2\kappa\sinh(n\kappa/2)
 =2\delta\lambda\cosh(n\kappa/2),
\]
which is exactly \eqref{eq:telegraph-defect}.

Equivalently, the polarization entry of the homogeneous telegraph
resolvent is
\[
 \frac1n\sum_{\ell\in\Z}
 \frac{\lambda}{\lambda(\lambda+2\omega)+(2\pi\ell/n)^2}
 =\frac{\lambda}{2\kappa}\coth\left(\frac{n\kappa}{2}\right).
\]
The series converges normally away from the homogeneous spectrum and
the last expression supplies its meromorphic continuation when
$\kappa$ is purely imaginary, in agreement with the direct jump
calculation.  Finally, reflection-even modes have $m(0)=0$, so the jump
vanishes and those modes are protected exactly as in Lemma
\ref{lem:protected-sector}.
\end{proof}

\section{Conclusions}

\begin{remark}[The theorem package]\label{rem:package}
For the generator in Definition \ref{def:generator}:
\begin{enumerate}
\item Theorem \ref{thm:characteristic} and Theorem \ref{thm:transfer}
give exact characteristic equations for every $n$, $\omega$, and
$\omega_0$; the trace equation is used away from its denominators and
its cleared polynomial identity is used at their zeros.
\item Theorems \ref{thm:jordan-regular} and \ref{thm:flat-band}
classify the algebraic and geometric multiplicities of every eigenvalue,
including every exceptional point.
\item Theorem \ref{thm:localization} gives an exact criterion and decay
length for a defect-localized mode.
\item Theorems \ref{thm:fixed-gap} and \ref{thm:critical-gap} give the
fixed-rate, vanishing-defect, and bulk-critical large-$n$ gap
asymptotics.
\item Theorem \ref{thm:telegraph-comparison} and Proposition
\ref{prop:telegraph-defect} compare the exact relaxation rate with the
standard and diffusion-corrected telegraph approximations.
\end{enumerate}
Each item retains the hypotheses stated in its cited theorem: in
particular, localization is asserted off the bulk spectrum (with the
compact case treated separately), finite-ring localization convergence
assumes a simple noncompact root, and the gap results use their stated
fixed-rate or double-scaling regimes.
\end{remark}

\section*{AI-use disclosure}

The proof architecture and manuscript were developed with assistance from
OpenAI GPT-5.6 Sol (\texttt{gpt-5.6-sol}) at max reasoning effort.  The model
served as an interactive research assistant for deriving and testing
arguments, checking citations and calculations, writing exact verification
code, and editing the exposition.  The author set the research direction,
evaluated and approved the final arguments, and assumes responsibility for
the mathematical claims, references, computations, and text.

\bibliographystyle{alpha}
\bibliography{refs}

\end{document}
